% Options for packages loaded elsewhere
% Options for packages loaded elsewhere
\PassOptionsToPackage{unicode}{hyperref}
\PassOptionsToPackage{hyphens}{url}
\PassOptionsToPackage{dvipsnames,svgnames,x11names}{xcolor}
%
\documentclass[
  letterpaper,
  DIV=11,
  numbers=noendperiod]{scrartcl}
\usepackage{xcolor}
\usepackage{amsmath,amssymb}
\setcounter{secnumdepth}{-\maxdimen} % remove section numbering
\usepackage{iftex}
\ifPDFTeX
  \usepackage[T1]{fontenc}
  \usepackage[utf8]{inputenc}
  \usepackage{textcomp} % provide euro and other symbols
\else % if luatex or xetex
  \usepackage{unicode-math} % this also loads fontspec
  \defaultfontfeatures{Scale=MatchLowercase}
  \defaultfontfeatures[\rmfamily]{Ligatures=TeX,Scale=1}
\fi
\usepackage{lmodern}
\ifPDFTeX\else
  % xetex/luatex font selection
\fi
% Use upquote if available, for straight quotes in verbatim environments
\IfFileExists{upquote.sty}{\usepackage{upquote}}{}
\IfFileExists{microtype.sty}{% use microtype if available
  \usepackage[]{microtype}
  \UseMicrotypeSet[protrusion]{basicmath} % disable protrusion for tt fonts
}{}
\makeatletter
\@ifundefined{KOMAClassName}{% if non-KOMA class
  \IfFileExists{parskip.sty}{%
    \usepackage{parskip}
  }{% else
    \setlength{\parindent}{0pt}
    \setlength{\parskip}{6pt plus 2pt minus 1pt}}
}{% if KOMA class
  \KOMAoptions{parskip=half}}
\makeatother
% Make \paragraph and \subparagraph free-standing
\makeatletter
\ifx\paragraph\undefined\else
  \let\oldparagraph\paragraph
  \renewcommand{\paragraph}{
    \@ifstar
      \xxxParagraphStar
      \xxxParagraphNoStar
  }
  \newcommand{\xxxParagraphStar}[1]{\oldparagraph*{#1}\mbox{}}
  \newcommand{\xxxParagraphNoStar}[1]{\oldparagraph{#1}\mbox{}}
\fi
\ifx\subparagraph\undefined\else
  \let\oldsubparagraph\subparagraph
  \renewcommand{\subparagraph}{
    \@ifstar
      \xxxSubParagraphStar
      \xxxSubParagraphNoStar
  }
  \newcommand{\xxxSubParagraphStar}[1]{\oldsubparagraph*{#1}\mbox{}}
  \newcommand{\xxxSubParagraphNoStar}[1]{\oldsubparagraph{#1}\mbox{}}
\fi
\makeatother


\usepackage{longtable,booktabs,array}
\newcounter{none} % for unnumbered tables
\usepackage{calc} % for calculating minipage widths
% Correct order of tables after \paragraph or \subparagraph
\usepackage{etoolbox}
\makeatletter
\patchcmd\longtable{\par}{\if@noskipsec\mbox{}\fi\par}{}{}
\makeatother
% Allow footnotes in longtable head/foot
\IfFileExists{footnotehyper.sty}{\usepackage{footnotehyper}}{\usepackage{footnote}}
\makesavenoteenv{longtable}
\usepackage{graphicx}
\makeatletter
\newsavebox\pandoc@box
\newcommand*\pandocbounded[1]{% scales image to fit in text height/width
  \sbox\pandoc@box{#1}%
  \Gscale@div\@tempa{\textheight}{\dimexpr\ht\pandoc@box+\dp\pandoc@box\relax}%
  \Gscale@div\@tempb{\linewidth}{\wd\pandoc@box}%
  \ifdim\@tempb\p@<\@tempa\p@\let\@tempa\@tempb\fi% select the smaller of both
  \ifdim\@tempa\p@<\p@\scalebox{\@tempa}{\usebox\pandoc@box}%
  \else\usebox{\pandoc@box}%
  \fi%
}
% Set default figure placement to htbp
\def\fps@figure{htbp}
\makeatother





\setlength{\emergencystretch}{3em} % prevent overfull lines

\providecommand{\tightlist}{%
  \setlength{\itemsep}{0pt}\setlength{\parskip}{0pt}}



 


\KOMAoption{captions}{tableheading}
\makeatletter
\@ifpackageloaded{caption}{}{\usepackage{caption}}
\AtBeginDocument{%
\ifdefined\contentsname
  \renewcommand*\contentsname{Table of contents}
\else
  \newcommand\contentsname{Table of contents}
\fi
\ifdefined\listfigurename
  \renewcommand*\listfigurename{List of Figures}
\else
  \newcommand\listfigurename{List of Figures}
\fi
\ifdefined\listtablename
  \renewcommand*\listtablename{List of Tables}
\else
  \newcommand\listtablename{List of Tables}
\fi
\ifdefined\figurename
  \renewcommand*\figurename{Figure}
\else
  \newcommand\figurename{Figure}
\fi
\ifdefined\tablename
  \renewcommand*\tablename{Table}
\else
  \newcommand\tablename{Table}
\fi
}
\@ifpackageloaded{float}{}{\usepackage{float}}
\floatstyle{ruled}
\@ifundefined{c@chapter}{\newfloat{codelisting}{h}{lop}}{\newfloat{codelisting}{h}{lop}[chapter]}
\floatname{codelisting}{Listing}
\newcommand*\listoflistings{\listof{codelisting}{List of Listings}}
\makeatother
\makeatletter
\makeatother
\makeatletter
\@ifpackageloaded{caption}{}{\usepackage{caption}}
\@ifpackageloaded{subcaption}{}{\usepackage{subcaption}}
\makeatother
\usepackage{bookmark}
\IfFileExists{xurl.sty}{\usepackage{xurl}}{} % add URL line breaks if available
\urlstyle{same}
\hypersetup{
  pdftitle={Ujian Tengah Semester Probabilitas dan Statistika 2022/2023},
  colorlinks=true,
  linkcolor={blue},
  filecolor={Maroon},
  citecolor={Blue},
  urlcolor={Blue},
  pdfcreator={LaTeX via pandoc}}


\title{Ujian Tengah Semester Probabilitas dan Statistika 2022/2023}
\usepackage{etoolbox}
\makeatletter
\providecommand{\subtitle}[1]{% add subtitle to \maketitle
  \apptocmd{\@title}{\par {\large #1 \par}}{}{}
}
\makeatother
\subtitle{Lembar Soal dan Solusi (Asumsi 3 Digit Terakhir NIM: 123)}
\author{}
\date{}
\begin{document}
\maketitle

\renewcommand*\contentsname{Table of contents}
{
\setcounter{tocdepth}{3}
\tableofcontents
}

\section{LEMBAR SOAL}\label{lembar-soal}

\textbf{Peraturan:} 1. Waktu untuk mengerjakan ujian adalah 120 menit.
2. Sifat ujian adalah \emph{open book, open notes and tools} (python
notebook, excel, R, calculator, dll). 3. Dilarang membuka aplikasi yang
dapat digunakan untuk \emph{chatting} selama ujian. 4. Dilarang
berkomunikasi dan bekerja sama dengan orang lain.

\begin{center}\rule{0.5\linewidth}{0.5pt}\end{center}

\subsubsection{\texorpdfstring{\textbf{Soal 1: Rapid Test CoVID-19 (Skor
25)}}{Soal 1: Rapid Test CoVID-19 (Skor 25)}}\label{soal-1-rapid-test-covid-19-skor-25}

Diketahui bahwa kasus CoVID-19 aktif di Indonesia pada bulan Oktober
2022 berjumlah \textbf{17.423} orang. Populasi di Indonesia diketahui
berjumlah 270.600.000 jiwa. Prevalensi adalah proporsi jumlah kasus
aktif terhadap populasi.

Sebuah rapid test digunakan untuk mendeteksi infeksi CoVID-19. Hasil
evaluasi rapid test yang digunakan adalah sebagai berikut: \emph{False
Positive Rate} (FPR) dari rapid test ini adalah \textbf{0,0003}.
\emph{False Negative Rate} (FNR) dari rapid test ini adalah
\textbf{0,203}.

Tuliskan jawaban Anda dalam bilangan desimal (tanpa persen): a. Nilai
prevalensi kasus CoVID-19 di Indonesia adalah\ldots{} (bulatkan menjadi
5 angka di belakang koma desimal) b. Probabilitas hasil rapid test
reaktif apabila dipastikan terinfeksi CoVID-19 adalah\ldots{} (bulatkan
menjadi 3 angka desimal) c.~Probabilitas hasil rapid test TIDAK reaktif
apabila dipastikan terinfeksi CoVID-19 adalah\ldots{} (bulatkan menjadi
3 angka desimal) d.~Probabilitas hasil rapid test reaktif apabila
dipastikan TIDAK terinfeksi CoVID-19 adalah\ldots{} (bulatkan menjadi 4
angka desimal) e. Probabilitas hasil rapid test TIDAK reaktif apabila
dipastikan TIDAK terinfeksi CoVID-19 adalah\ldots{} (bulatkan menjadi 4
angka desimal) f.~Probabilitas seseorang terinfeksi CoVID-19 apabila
dipastikan hasil rapid test nya reaktif adalah\ldots{} (bulatkan menjadi
4 angka desimal) g. Probabilitas seseorang TIDAK terinfeksi CoVID-19
apabila dipastikan hasil rapid test nya reaktif adalah\ldots{} (bulatkan
menjadi 4 angka desimal) h. Probabilitas seseorang terinfeksi CoVID-19
apabila dipastikan hasil rapid test nya TIDAK reaktif adalah\ldots{}
(bulatkan menjadi 6 angka desimal) i. Probabilitas seseorang TIDAK
terinfeksi CoVID-19 apabila dipastikan hasil rapid test nya TIDAK
reaktif adalah\ldots{} (bulatkan menjadi 6 angka desimal) j.
Probabilitas hasil rapid test reaktif saat seseorang di Indonesia diuji
adalah\ldots{} (bulatkan menjadi 5 angka desimal) k. Probabilitas hasil
rapid test TIDAK reaktif saat seseorang di Indonesia diuji
adalah\ldots{} (bulatkan menjadi 5 angka desimal) l. Akurasi dari rapid
test tersebut adalah\ldots{} (bulatkan menjadi 5 angka desimal) m.
Diadakan suatu festival yang dihadiri oleh \textbf{1.023} penduduk
Indonesia. Asumsikan status independen. Probabilitas minimal 1 orang
yang menghadiri festival tersebut terinfeksi CoVID-19 adalah\ldots{}
(bulatkan menjadi 3 angka desimal) n.~Diketahui nilai \emph{Case
Fatality Rate} CoVID-19 saat ini di Indonesia adalah 0,43\% dan saat ini
terdapat \textbf{17.423} orang yang terpapar. Probabilitas tidak ada
korban jiwa sama sekali dari para penderita yang terpapar ini
adalah\ldots{} (bulatkan menjadi 2 angka desimal)

\begin{center}\rule{0.5\linewidth}{0.5pt}\end{center}

\subsubsection{\texorpdfstring{\textbf{Soal 2: Fraud Detection System
(Skor
10)}}{Soal 2: Fraud Detection System (Skor 10)}}\label{soal-2-fraud-detection-system-skor-10}

Sebuah perusahaan digital \emph{startup} PT. P meletakkan aplikasi Fraud
Detection System (FDS) di \textbf{4 buah cloud} yang berbeda dan
independen. Masing-masing cloud memiliki probabilitas untuk bekerja
dengan baik sebesar \textbf{0,973}. Definisikan \(X\) sebagai random
variable yang merepresentasikan jumlah cloud PT. P yang bekerja dengan
baik.

\begin{enumerate}
\def\labelenumi{\alph{enumi}.}
\tightlist
\item
  Probabilitas jumlah cloud yang bekerja dengan baik tidak kurang dari 1
  adalah\ldots{} (bulatkan menjadi 2 angka desimal)
\item
  Probabilitas jumlah cloud yang bekerja dengan baik tidak kurang dari 2
  adalah\ldots{} (bulatkan menjadi 6 angka desimal)
\item
  Probabilitas jumlah cloud yang bekerja dengan baik tidak lebih dari 3
  adalah\ldots{} (bulatkan menjadi 4 angka desimal)
\item
  PT. P memiliki 8000 pelanggan. Apabila tiap cloud mampu memenuhi
  kebutuhan 3000 pelanggan, berapa probabilitas terdapat pelanggan PT. P
  yang kebutuhannya tidak terpenuhi oleh cloud? (bulatkan menjadi 4
  angka desimal)
\item
  Nilai ekspektasi dari random variable \(X\) adalah\ldots{} (bulatkan
  menjadi 3 angka desimal)
\item
  Nilai variansi dari random variable \(X\) adalah\ldots{} (bulatkan
  menjadi 3 angka desimal)
\end{enumerate}

\begin{center}\rule{0.5\linewidth}{0.5pt}\end{center}

\subsubsection{\texorpdfstring{\textbf{Soal 3: Tegangan dan Daya (Skor
10)}}{Soal 3: Tegangan dan Daya (Skor 10)}}\label{soal-3-tegangan-dan-daya-skor-10}

Tegangan dari suatu tahanan 40 Ohm merupakan suatu random variable yang
terdistribusi normal dengan mean \textbf{32} Volt dan standar deviasi
\textbf{5} Volt. Daya \(P\) yang didisipasikan dalam tahanan tersebut
merupakan suatu random variable juga (\(P = \frac{V^2}{R}\)).

\begin{enumerate}
\def\labelenumi{\alph{enumi}.}
\tightlist
\item
  Probabilitas tegangan yang melalui tahanan tersebut ada di antara 20
  -- 25 Volt adalah\ldots{} (bulatkan menjadi 4 angka desimal)
\item
  Probabilitas tegangan yang melalui tahanan tersebut \textgreater{} 30
  Volt adalah\ldots{} (bulatkan menjadi 4 angka desimal)
\item
  Probabilitas tegangan yang melalui tahanan tersebut \textless{} 25
  Volt adalah\ldots{} (bulatkan menjadi 4 angka desimal)
\item
  Nilai ekspektasi dari daya yang terdisipasi pada tegangan tersebut
  adalah\ldots{} Watt (bulatkan menjadi 3 angka desimal)
\end{enumerate}

\begin{center}\rule{0.5\linewidth}{0.5pt}\end{center}

\subsubsection{\texorpdfstring{\textbf{Soal 4: Transmisi Paket (Skor
15)}}{Soal 4: Transmisi Paket (Skor 15)}}\label{soal-4-transmisi-paket-skor-15}

Sebuah paket sepanjang \textbf{256 bit} ditransmisikan melalui suatu
kanal ber-noise dengan mekanisme koreksi error yang dapat mengoreksi
\textbf{tidak lebih dari 12 error}. Probabilitas bahwa sebuah bit error
adalah \textbf{0,032} untuk semua bit secara independen.

\begin{enumerate}
\def\labelenumi{\alph{enumi}.}
\tightlist
\item
  Probabilitas bahwa paket tersebut akan diterima dengan benar
  adalah\ldots{} (bulatkan menjadi 4 angka desimal)
\item
  Nilai ekspektasi dari jumlah error adalah\ldots{} (bulatkan menjadi 2
  angka desimal)
\item
  Nilai variansi dari jumlah error adalah\ldots{} (bulatkan menjadi 2
  angka desimal)
\item
  \textbf{253} buah paket sepanjang 256 bit ditransmisikan. Probabilitas
  terdapat tidak kurang dari 3 paket ditransmisikan dalam kondisi error
  adalah\ldots{} (bulatkan menjadi 8 angka desimal)
\item
  \textbf{253} buah paket ditransmisikan. Variansi dari jumlah paket
  yang ditransmisikan dalam kondisi error adalah\ldots{} (bulatkan
  menjadi 3 angka desimal)
\end{enumerate}

\begin{center}\rule{0.5\linewidth}{0.5pt}\end{center}

\subsubsection{\texorpdfstring{\textbf{Soal 5: Sistem Komunikasi (Skor
20)}}{Soal 5: Sistem Komunikasi (Skor 20)}}\label{soal-5-sistem-komunikasi-skor-20}

\emph{(Asumsikan terdapat gambar graf topologi yang mendeskripsikan 11
link, i hingga xi, dan node A-J)}. Dalam kebijakan pemeliharaan saat
ini, probabilitas dari sebuah link bekerja dengan baik bernilai
\textbf{0,93}.

\begin{enumerate}
\def\labelenumi{\alph{enumi}.}
\tightlist
\item
  Diketahui tepat 7 link komunikasi tidak bekerja (4 link dipastikan
  bekerja), probabilitas D dapat berkomunikasi dengan J adalah\ldots{}
\item
  Diketahui tepat 8 link komunikasi tidak bekerja (3 link dipastikan
  bekerja), probabilitas D dapat berkomunikasi dengan J adalah\ldots{}
\item
  Diketahui link i, v, vii, ix dan x dipastikan tidak bekerja,
  probabilitas I dapat berkomunikasi dengan J adalah\ldots{}
\item
  Diketahui link ii, iv, vi, viii, x, xi dipastikan tidak bekerja,
  probabilitas B dapat berkomunikasi dengan G adalah\ldots{}
\end{enumerate}

\begin{center}\rule{0.5\linewidth}{0.5pt}\end{center}

\subsubsection{\texorpdfstring{\textbf{Soal 6: Tender Pengadaan Server
(Skor
20)}}{Soal 6: Tender Pengadaan Server (Skor 20)}}\label{soal-6-tender-pengadaan-server-skor-20}

Diketahui PT. G dan PT. H bersaing memenangkan tender. * Rumus nilai
akhir = 65\% nilai teknis + 35\% nilai harga. * Rumus nilai harga = (50
M -- Harga diajukan)/50 M * 100. * Nilai teknis PT. G adalah \textbf{63}
dan PT. H adalah \textbf{73}. * Harga PT. G terdistribusi seragam
(uniform) antara 0-30 M. * Harga PT. H terdistribusi seragam (uniform)
antara 0-45 M.

\begin{enumerate}
\def\labelenumi{\alph{enumi}.}
\tightlist
\item
  Nilai ekspektasi harga yang diajukan PT. G adalah \ldots{} milyar (2
  desimal)
\item
  Nilai variansi harga yang diajukan PT. H adalah \ldots{} milyar (2
  desimal)
\item
  Apabila \(X\) mewakili harga PT. G dan \(Y\) harga PT. H, berapa nilai
  \(f(x,y)\)? (5 desimal)
\item
  Probabilitas PT. G memenangkan tender adalah\ldots{} (3 desimal)
\item
  Probabilitas PT. H memenangkan tender adalah\ldots{} (3 desimal)
\end{enumerate}

\newpage

\section{LEMBAR JAWABAN / SOLUSI}\label{lembar-jawaban-solusi}

\subsubsection{\texorpdfstring{\textbf{Solusi Soal
1}}{Solusi Soal 1}}\label{solusi-soal-1}

Berdasarkan parameter dari substitusi NIM 1,2,3: * Prevalensi = Peluang
terinfeksi \(\approx P(D) = 17423 / 270.600.000 = 0,000064386...\) *
Peluang tidak terinfeksi \(\approx P(D^c) = 1 - P(D) = 0,9999356...\) *
False Positive Rate (FPR)
\(= P(+|D^c) = 0,0003 \implies P(-|D^c) = 0,9997\) * False Negative Rate
(FNR) \(= P(-|D) = 0,203 \implies P(+|D) = 0,797\)

\textbf{Jawaban:} \textbf{a.}
\(P(D) = \frac{17423}{270600000} \approx \mathbf{0,00006}\) \textbf{b.}
\(P(+|D) = 1 - FNR = 1 - 0,203 = \mathbf{0,797}\) \textbf{c.}
\(P(-|D) = FNR = \mathbf{0,203}\) \textbf{d.}
\(P(+|D^c) = FPR = \mathbf{0,0003}\) \textbf{e.}
\(P(-|D^c) = 1 - FPR = \mathbf{0,9997}\)

(Langkah Perantara untuk Probabilitas Total):
\(P(+) = P(+|D)P(D) + P(+|D^c)P(D^c) = (0,797)(0,0000644) + (0,0003)(0,9999356) = 0,0000513 + 0,0003 = 0,0003513\)
\(P(-) = 1 - P(+) = 0,9996487\)

\textbf{f.}
\(P(D|+) = \frac{P(+|D)P(D)}{P(+)} = \frac{0,0000513}{0,0003513} \approx \mathbf{0,1461}\)
\textbf{g.} \(P(D^c|+) = 1 - P(D|+) = \mathbf{0,8539}\) \textbf{h.}
\(P(D|-) = \frac{P(-|D)P(D)}{P(-)} = \frac{(0,203)(0,0000644)}{0,9996487} = \frac{0,00001307}{0,9996487} \approx \mathbf{0,000013}\)
\textbf{i.} \(P(D^c|-) = 1 - P(D|-) \approx \mathbf{0,999987}\)
\textbf{j.} \(P(+) \approx \mathbf{0,00035}\) \textbf{k.}
\(P(-) \approx \mathbf{0,99965}\) \textbf{l.}
\(\text{Akurasi} = P(+|D)P(D) + P(-|D^c)P(D^c) = 0,0000513 + 0,9996356 \approx \mathbf{0,99969}\)
\textbf{m.} Distribusi binomial dengan \(n=1023, p=0,0000644\):
\(P(\text{Kasus} \ge 1) = 1 - (1-p)^{1023} = 1 - 0,93627 \approx \mathbf{0,064}\)
\textbf{n.}
\(P(\text{0 Jiwa}) = (1 - 0,0043)^{17423} = (0,9957)^{17423} \approx e^{-17423 \times 0,0043} = e^{-74,91} \approx \mathbf{0,00}\)

\begin{center}\rule{0.5\linewidth}{0.5pt}\end{center}

\subsubsection{\texorpdfstring{\textbf{Solusi Soal
2}}{Solusi Soal 2}}\label{solusi-soal-2}

Distribusi Binomial \(X \sim Bin(n=4, p=0,973)\), dengan \(q = 0,027\).
\textbf{a.}
\(P(X \ge 1) = 1 - P(X=0) = 1 - (0,027)^4 \approx \mathbf{1,00}\)
\textbf{b.}
\(P(X \ge 2) = 1 - P(X=0) - P(X=1) = 1 - (0,027)^4 - 4(0,973)(0,027)^3 \approx 1 - 0,0000766 = \mathbf{0,999923}\)
\textbf{c.}
\(P(X \le 3) = 1 - P(X=4) = 1 - (0,973)^4 = 1 - 0,8963 = \mathbf{0,1037}\)
\textbf{d.} Pelanggan tidak terpenuhi jika cloud yang hidup tidak mampu
menampung (3 cloud \(\times\) 3000 = 9000). Karena 8000 butuh 3 cloud,
maka gagal terpenuhi jika cloud hidup \(X < 3\) atau \(X \le 2\).
\(P(X \le 2) = P(X=0) + P(X=1) + P(X=2) = 0,0000005 + 0,0000766 + 6(0,973)^2(0,027)^2 = \mathbf{0,0042}\)
\textbf{e.} Ekspektasi \(E(X) = np = 4(0,973) = \mathbf{3,892}\)
\textbf{f.} Variansi \(Var(X) = npq = 4(0,973)(0,027) = \mathbf{0,105}\)

\begin{center}\rule{0.5\linewidth}{0.5pt}\end{center}

\subsubsection{\texorpdfstring{\textbf{Solusi Soal
3}}{Solusi Soal 3}}\label{solusi-soal-3}

\(V \sim N(32, 5^2)\). Transformasi Z: \(Z = \frac{V - 32}{5}\)
\textbf{a.}
\(P(20 \le V \le 25) = P(-2,4 \le Z \le -1,4) = P(Z < -1,4) - P(Z < -2,4) = 0,0808 - 0,0082 = \mathbf{0,0726}\)
\textbf{b.} \(P(V > 30) = P(Z > -0,4) = 1 - 0,3446 = \mathbf{0,6554}\)
\textbf{c.} \(P(V < 25) = P(Z < -1,4) = \mathbf{0,0808}\) \textbf{d.}
\(E(P) = E(V^2 / 40) = \frac{1}{40} E(V^2)\). Sifat variansi:
\(E(V^2) = Var(V) + (E(V))^2 = 25 + (32)^2 = 25 + 1024 = 1049\).
\(E(P) = 1049 / 40 = \mathbf{26,225}\)

\begin{center}\rule{0.5\linewidth}{0.5pt}\end{center}

\subsubsection{\texorpdfstring{\textbf{Solusi Soal
4}}{Solusi Soal 4}}\label{solusi-soal-4}

\(Y \sim Bin(256, 0,032)\). \textbf{a.} \(P(Y \le 12)\). (Dapat
diestimasi menggunakan pendekatan normal karena N cukup besar).
\(\mu = np = 8,192\),
\(\sigma = \sqrt{npq} = \sqrt{7,929} \approx 2,816\).
\(Z = \frac{12,5 - 8,192}{2,816} \approx 1,53\).
\(P(Z < 1,53) \approx \mathbf{0,9370}\). \textbf{b.}
\(E(Y) = np = 256(0,032) = \mathbf{8,19}\) \textbf{c.}
\(Var(Y) = npq = 256(0,032)(0,968) = \mathbf{7,93}\) \textbf{d.}
Terdapat \(M=253\) paket. Peluang satu paket gagal adalah
\(p_{gagal} = 1 - 0,9370 = 0,063\). Variabel paket gagal
\(W \sim Bin(253, 0,063)\).
\(P(W \ge 3) = 1 - P(W \le 2) = 1 - \sum_{k=0}^{2} \binom{253}{k}(0,063)^k(0,937)^{253-k}\).
\emph{(Nilainya mendekati \(\approx \mathbf{0,99912061}\) jika dihitung
eksak dengan kalkulator).} \textbf{e.}
\(Var(W) = M \cdot p_{gagal} \cdot (1 - p_{gagal}) = 253(0,063)(0,937) \approx \mathbf{14,932}\)

\begin{center}\rule{0.5\linewidth}{0.5pt}\end{center}

\subsubsection{\texorpdfstring{\textbf{Solusi Soal
5}}{Solusi Soal 5}}\label{solusi-soal-5}

\emph{(Soal ini memerlukan gambar graf topologi jaringan spesifik untuk
menentukan enumerasi / path matrix. Solusi ini mendeskripsikan secara
matematis jika model diasumsikan telah ada)}. Jika diasumsikan kita
memiliki jumlah tautan (edges) dan rute, pendekatannya adalah sebagai
berikut: \textbf{a \& b.} Jika diketahui tepat \(n\) link tidak bekerja,
maka kombinasi penentuan ``sukses'' murni didasarkan pada permutasi
subgraf.
\(P(\text{sukses}) = \frac{\text{Jumlah kombinasi set link sukses dari } C(\text{total link}, \text{link aktif})}{\text{Total Kombinasi}}\).
\textbf{c \& d.} Apabila ada \(K\) set link yang ``dipastikan mati'',
maka edge (garis) tersebut dapat dihilangkan secara langsung dari
topologi \emph{Network Flow}. Setelah grafik direduksi, peluang
komunikasi dihitung melalui dekomposisi \emph{Reliability Block Diagram}
(Seri/Paralel) pada tautan yang tersisa yang masing-masing berpeluang
hidup \(p=0,93\).

\begin{center}\rule{0.5\linewidth}{0.5pt}\end{center}

\subsubsection{\texorpdfstring{\textbf{Solusi Soal
6}}{Solusi Soal 6}}\label{solusi-soal-6}

Harga G \(= X \sim Uniform(0, 30)\). Harga H
\(= Y \sim Uniform(0, 45)\). Formula Final:
\(\text{NA}_G = 0,65(63) + 0,35 \left( \frac{50-X}{50} \times 100 \right) = 40,95 + 35 - 0,7X = 75,95 - 0,7X\)
\(\text{NA}_H = 0,65(73) + 0,35 \left( \frac{50-Y}{50} \times 100 \right) = 47,45 + 35 - 0,7Y = 82,45 - 0,7Y\)
G menang jika
\(\text{NA}_G > \text{NA}_H \implies 75,95 - 0,7X > 82,45 - 0,7Y \implies 0,7Y - 0,7X > 6,5 \implies Y - X > 9,2857\).

\textbf{Jawaban:} \textbf{a.}
\(E(X) = \frac{0 + 30}{2} = \mathbf{15,00}\) M \textbf{b.}
\(Var(Y) = \frac{(45 - 0)^2}{12} = \frac{2025}{12} = \mathbf{168,75}\) M
\textbf{c.} \(f(x,y)\) untuk distribusi gabungan yang independen
\(= f(x) \cdot f(y) = \frac{1}{30} \times \frac{1}{45} = \frac{1}{1350} \approx \mathbf{0,00074}\)
\textbf{d.} Probabilitas G menang (\(P(Y - X > 9,2857)\)): Menggunakan
integral geometri pada kotak \$ \times \$ untuk area \(Y > X + 9,2857\):
Luas Area \(= \int_{0}^{30} \int_{x+9,2857}^{45} \frac{1}{1350} dy dx\)
Luas
\(= \frac{1}{1350} \left[ \int_{0}^{30} (35,7143 - x) dx \right] = \frac{1}{1350} \left[ 35,7143(30) - \frac{900}{2} \right] = \frac{1071,429 - 450}{1350} = \frac{621,429}{1350} \approx \mathbf{0,460}\)
\textbf{e.} Probabilitas H menang
\(= 1 - P(G \text{ menang}) = 1 - 0,460 = \mathbf{0,540}\)




\end{document}
